\documentclass{my_memo}

\usepackage{booktabs}
\usepackage{tabularx}
\usepackage{array}
\usepackage{listings}
\sloppy

% The my_memo class sits close to the 16 math-family limit (mathrsfs +
% amssymb + kotex + explicit \DeclareMathAlphabet for \mathbf/\mathit/
% \mathbit).  Route \mathbf through \boldsymbol so no new family is
% allocated.
\let\mathbf\boldsymbol

\lstset{
  language=Fortran,
  basicstyle=\ttfamily\small,
  keywordstyle=\color{blue!60!black},
  commentstyle=\color{green!40!black},
  stringstyle=\color{red!60!black},
  columns=fullflexible,
  breaklines=true,
  frame=single,
  framesep=4pt,
  showstringspaces=false,
  numbers=none
}
\lstdefinestyle{cpp}{language=C++,keywordstyle=\color{blue!60!black}}
\lstdefinestyle{sh}{language=bash,keywordstyle=\color{black},basicstyle=\ttfamily\small}

\newcommand{\code}[1]{\texttt{#1}}
\newcommand{\um}{\ensuremath{\,\mu\mathrm{m}}}
\newcommand{\Cext}{\ensuremath{C_{\rm ext}}}
\newcommand{\Cabs}{\ensuremath{C_{\rm abs}}}
\newcommand{\Cpol}{\ensuremath{C_{\rm pol}}}
\newcommand{\Ccirc}{\ensuremath{C_{\rm circ}}}
\newcommand{\Qabs}{\ensuremath{Q_{\rm abs}}}
\newcommand{\lamIpol}{\ensuremath{(\lambda I_\lambda)_{\rm pol}}}
\newcommand{\aalign}{\ensuremath{a_{\rm align}}}
\newcommand{\adisr}{\ensuremath{a_{\rm disr}}}
\newcommand{\falign}{\ensuremath{f_{\rm align}}}
\newcommand{\fhighJ}{\ensuremath{f_{\rm high\text{-}J}}}
\newcommand{\Av}{\ensuremath{A_V}}

\title{Polarization by Aligned Grains: What SEDust and MoCafe Would Need}
\author{Kwang-il Seon}
\date{2026 July 18}

\begin{document}
\maketitle

\begin{abstract}
\noindent
This memo audits what it would take to replace the empirical dichroic
polarization treatment of Seon (2018, ApJ 862, 87) with genuine
aligned-grain physics in the current SEDust\,$+$\,MoCafe stack. The 2018
model is dissected first: its \emph{scattering} polarization was already
exact (full numerical Mie Mueller matrix with Stokes tracing), but its
\emph{dichroic} polarization was an add-on fitting formula with no
alignment physics, no wavelength dependence, and additive rather than
multiplicative accumulation of Stokes $Q$ and $U$. The three ingredients
that would be needed instead are then examined in turn: polarized cross
sections, an alignment prescription, and a polarized transfer solver. The
central practical finding is that the first of these is already in hand.
The file \code{q\_DH21Ad\_P0.20\_Fe0.00\_1.400.dat.gz} stores three
separate orientation blocks rather than an orientation average, so
$\Cpol$ can be built today without any new T-matrix computation; values
derived from it reproduce the Hensley \& Draine (2023) released
\code{polarized\_extinction.dat} to four significant figures. A staged
plan follows, in which stage~1 (SEDust cross sections and an emission
accumulator) is a few days of low-risk work, stage~2 already permits a
quantitative comparison against the SOFIA/HAWC+ 154\um\ maps of NGC~891,
and only stage~4 (self-consistent alignment) requires the invasive change
of a vectorial radiation-field pass. Validation targets and the points
that remain unverified are listed explicitly.
\end{abstract}

\tableofcontents

%====================================================================
\section{Motivation: anatomy of the 2018 model}
\label{sec:motivation}

The question this memo answers is narrow: if Seon (2018, ApJ 862, 87;
\code{2018ApJ...862...87S}), ``Polarization as a Probe of Thick Dust
Disks in Edge-on Galaxies: Application to NGC~891'', were redone now,
what would have to be built? Answering it requires an exact statement of
what that model did and did not contain. The description below is taken
from arXiv:1806.06101v1.

\subsection{What was already correct: scattering polarization}

The scattering half of the 2018 calculation was not an approximation.
Spherical grains were treated with the full numerical Mie Mueller matrix,
retaining all four independent elements $S_{11}$, $S_{12}$, $S_{33}$,
$S_{34}$; photons carried a Stokes vector; the scattering and peel-off
bookkeeping followed the algorithm of Peest et al.\ (2017, A\&A 601,
A92; \code{2017A\&A...601A..92P}); and the implementation was verified
against the polarized scattering benchmark of Bianchi, Ferrara \& Giovanardi
(1996, ApJ 465, 127; \code{1996ApJ...465..127B}). No Henyey--Greenstein
or Rayleigh proxy was used for the phase matrix. That part of the code
needs no revision.

\subsection{What was empirical: dichroic polarization}

The dichroic half descends from the prescriptions of Wood (1997, ApJ 477,
L25; \code{1997ApJ...477L..25W}) and Wood \& Jones (1997, AJ 114, 1405;
\code{1997AJ....114.1405W}). Seon (2018) adopted the Whittet et al.\
(2008, ApJ 674, 304; \code{2008ApJ...674..304W}) empirical relation
$P_K/\tau_K \approx 5.08\,\tau_V^{-0.52}$ and converted it to
\begin{equation}
  P_V \;=\; 2.7\,\tau_V^{0.48} \quad [\%],
  \label{eq:whittet}
\end{equation}
then, at every interaction site, added a polarized fraction
$P_V\sin^2\theta$ oriented parallel to the local magnetic field
projected on the sky, with $\theta$ the angle between the field and the
photon direction. The author's own summary in that paper is unambiguous:
``we assumed spherical dust grains when calculating the Mueller matrix
and adopted an empirical approach to take into account the dichroic
polarization effect.''

The internal tension is worth stating plainly. Under a strictly spherical
grain assumption the dichroic cross-section difference vanishes
identically, so the dichroic term was necessarily grafted on from
outside the optical model rather than derived from it.

\subsection{What was absent: alignment physics}

Radiative torque (RAT) alignment appears in the 2018 paper only as
introductory background. Neither the Rayleigh reduction factor $R$ nor an
alignment size threshold \aalign\ enters the calculation anywhere. All
alignment efficiency is absorbed into the single exponent $0.48$ and the
normalization of Eq.~(\ref{eq:whittet}), which are Milky Way averages.
That is a double-counting hazard as well as a physics gap: the Whittet
relation already folds in turbulent depolarization along real Galactic
sight lines, plus a mean alignment efficiency and a mean grain
composition, and the model then imposes its own explicit field geometry
on top of that same quantity.

The magnetic field was treated as two separate deterministic runs, purely
toroidal and purely vertical, compared against each other. There was no
turbulent component and no field-strength parameter. Crucially, the field
in that model sets only the polarization position angle and the $\sin^2\theta$
projection factor; it has no influence on alignment efficiency, because
no alignment efficiency is computed.

The dust model was MRN with Draine \& Lee (1984, ApJ 285, 89;
\code{1984ApJ...285...89D}) and Draine (2003, ApJ 598, 1017;
\code{2003ApJ...598.1017D}) silicate and graphite optical constants, with
PAHs excluded. The choice matters more than it looks: the maximum
single-scattering polarization is $36\%$ for that mixture but $24\%$ for
a WD01 mixture, a factor of $1.5$. Only the $V$ band was computed;
far-infrared and submillimeter polarized \emph{emission} was not a topic
of the paper at all.

\subsection{The implementation that survives in the tree}

The transport code was MoCafe with forced first scattering and
peeling-off. Dichroism was not applied by solving an extinction matrix;
polarization was simply incremented at each interaction site. Two
consequences follow directly: dichroic extinction never feeds back into
the optical depth, and birefringence (the $Q,U \leftrightarrow V$
conversion) is absent entirely.

The actual code is still present at
\code{\textasciitilde/MoCafe/Galaxy\_pol/dichroic.f90}. Lines 75--89 read:

\begin{lstlisting}
  select case(par%dichroic_method)
     case (2)
        pol = par%pol_dichroic * tau**(3.0_wp/4.0_wp)
     case (3)
        pol = par%pol_dichroic * tau**(0.48_wp)
     case (1)
        pol = par%pol_dichroic * tau**(0.48_wp) * sint**2
     case default
        pol = par%pol_dichroic * tau**(3.0_wp/4.0_wp) * sint**2
  end select
  T1    = cos2p * pol
  T2    = sin2p * pol
  photon%Q = photon%Q + T1
  photon%U = photon%U + T2
\end{lstlisting}

with the default amplitude \code{pol\_dichroic = 0.013} set in
\code{Galaxy\_pol/define.f90:210}. Two distinct physical defects are
visible here.

\paragraph{(i) The $\tau$ dependence is a fit, and the wavelength
dependence is missing.} The correct form for a uniform slab of aligned
grains is set by the polarized optical depth
$\Delta\tau = n_d \Cpol s$, giving
\begin{equation}
  p \;=\; \tanh(\Delta\tau),
  \label{eq:tanh}
\end{equation}
which reduces to $p \simeq \Delta\tau$ only in the optically thin limit.
Neither $\tau^{0.75}$ nor $\tau^{0.48}$ is that function. More seriously,
$\Cpol$ carries the entire wavelength dependence of dichroism, and it is
absent from the expression above: this formulation cannot reproduce a
Serkowski curve even in principle, because it has no $\lambda$ argument.

\paragraph{(ii) The accumulation is additive where the physics is
multiplicative.} The extinction matrix mixes $I$ into $Q$, so the correct
operation is a matrix exponential acting on the Stokes vector.
Incrementing $Q$ and $U$ by a fixed amount at each site agrees with that
only while $\tau \ll 1$; the two diverge once $\tau \gtrsim 1$, which is
precisely the regime a thick edge-on disk model is built to probe.

%====================================================================
\section{What has changed since 2018}
\label{sec:since2018}

\paragraph{Observations.} NGC~891 now has far-infrared polarimetry.
SOFIA/HAWC+ observed it at 154\um\ (Jones et al.\ 2020, AJ 160, 167;
\code{2020AJ....160..167J}), and a deeper analysis was published by Kim
et al.\ (2023, AJ 165, 223; \code{2023AJ....165..223K}). The conclusion
of the latter is directly in tension with the 2018 setup: the low
polarization fraction observed in the disk, together with its narrow
dispersion, requires many turbulent decorrelation cells along each line
of sight. A smooth, purely toroidal or purely vertical field cannot
produce that combination. Any redo therefore needs a turbulent field
component as a first-class ingredient, not an afterthought.

\paragraph{Methodology.} Two papers have since addressed exactly the
technical obstacles the 2018 implementation stepped around. Baes et al.\
(2019, A\&A 630, A61; \code{2019A\&A...630A..61B}) analyzed the
optical-depth problem in a dichroic medium, i.e.\ how to sample path
lengths when the extinction coefficient depends on the polarization state
being transported. Peest et al.\ (2023, A\&A 673, A112;
\code{2023A\&A...673A.112P}) provided a full 3D Monte Carlo radiative
transfer framework for aligned spheroidal grains, together with analytic
benchmarks. Both are directly usable.

%====================================================================
\section{Polarized optical properties}
\label{sec:optics}

\subsection{The minimum required set}

Three cross sections beyond the scalar ones are needed:
$\Cext^{\parallel}$ and $\Cext^{\perp}$ for dichroic extinction,
$\Cabs^{\parallel}$ and $\Cabs^{\perp}$ for polarized thermal emission,
and $\Ccirc$ for birefringence (the phase retardation that converts
linear into circular polarization). In practice these collapse into the
combinations $\Cext$, $\Cpol$, $\Ccirc$ defined below.

\subsection{Modified picket fence approximation}

The modified picket fence approximation (MPFA) of Draine \& Hensley
(2021b, ApJ 919, 65; \code{2021ApJ...919...65D}, Eqs.~11--17) reduces the
orientation dependence to two reference orientations. For an oblate
spheroid with symmetry axis $\hat a$, writing $\Theta$ for the angle
between the electric field and $\hat a$,
\begin{equation}
  C_E(\Theta) \;\simeq\; \cos^2\Theta\,C_E(0) + \sin^2\Theta\,C_E(90),
  \label{eq:mpfa}
\end{equation}
and identically for the magnetic-field-referenced $C_H$. Hensley \&
Draine (2023, ApJ 948, 55; hereafter HD23, \code{2023ApJ...948...55H})
then write the random-orientation and polarized cross sections as
\begin{align}
  C_{\rm ran} &= \tfrac13\left[C_E(0) + C_E(90) + C_H(90)\right],
    \label{eq:cran} \\
  C_{\rm pol}^{\rm oblate} &= \tfrac12\left[C_H(90) - C_E(90)\right],
    \label{eq:cpol}
\end{align}
their Eqs.~2 and 3 respectively. The accuracy of Eq.~(\ref{eq:mpfa}) is
excellent for $2\pi a/\lambda \lesssim 1$ and degrades to roughly $10\%$
near $a \sim \lambda$, but after integration over a realistic size
distribution the residual error is a few percent (DH21b).

\subsection{Partial alignment and the Rayleigh reduction factor}

Partial alignment enters through a single scalar, the Rayleigh reduction
factor
\begin{equation}
  R \;\equiv\; \tfrac32\langle\cos^2\beta\rangle - \tfrac12,
  \label{eq:Rdef}
\end{equation}
which is the same quantity HD23 call \falign. With $\theta$ the angle
between the magnetic field and the line of sight, the partially aligned
cross sections are
\begin{align}
  Q_{\rm abs}(\theta) &= \langle Q\rangle
     + R\,(Q_\parallel - Q_\perp)\left(\tfrac13 - \tfrac{\sin^2\theta}{2}\right),
     \label{eq:qabs} \\
  Q_{\rm abs,pol}(\theta) &= -\,R\,(Q_\parallel - Q_\perp)\,
     \tfrac{\sin^2\theta}{2}.
     \label{eq:qabspol}
\end{align}

The structural point for implementation is that the entire $\theta$
dependence factorizes into the single function $\sin^2\theta$. Only two
wavelength-by-size tables need to be stored, the orientation-averaged one
and the polarized one; the geometry factor is applied at run time in the
transport loop. No angular dimension is added to the tables.

The corresponding observables are given by HD23 Eq.~6 for polarized
extinction and HD23 Eq.~11 for polarized emission, the latter including
the integral $\int dT\,(dP/dT)$ over the stochastic temperature
distribution.

\subsection{Sign convention}

For an oblate spheroid $\hat a$ is the \emph{short} axis, and alignment
places $\hat a \parallel B$. Transmitted starlight is therefore polarized
parallel to the plane-of-sky field $B_{\rm sky}$, while thermal emission
is polarized perpendicular to it. Getting this backwards produces a
$90^\circ$ error that is easy to miss in an all-sky map and impossible to
miss in a comparison against HAWC+ vectors.

\subsection{The asset already in the tree}

The single most useful finding of this audit is that no new T-matrix
computation is required to get started. The file
\code{astrodust/data/dielectric/q\_DH21Ad\_P0.20\_Fe0.00\_1.400.dat.gz}
is not an orientation average. It contains three separate orientation
blocks for each of $Q_{\rm ext}$, $Q_{\rm abs}$, $Q_{\rm sca}$, on the
native $169$ sizes $\times$ $1129$ wavelengths grid. Its header states
verbatim:

\begin{lstlisting}[style=sh]
3 orientations: jori=1: k || a;; jori=2: k perp a, E || a;
                jori=3: k perp a, E perp a
\end{lstlisting}

so Eqs.~(\ref{eq:cran}) and (\ref{eq:cpol}) can be evaluated directly
from it.

This has been verified numerically, not merely asserted. Values built
from this table reproduce the HD23 released
\code{polarized\_extinction.dat} at $0.55$\um\ to four significant
figures (ratio $1.0000$), and a fit of the alignment function to the same
data returns $\aalign = 0.07492$\um, $\alpha = 1.799$, $f_{\max} = 1.0$,
identical to HD23 Table~1.

\subsection{Index convention traps}

The orientation index ordering is not consistent across Draine's own
distributions, and this is a live source of silent sign errors. Draine \&
Fraisse (2009, ApJ 696, 1; \code{2009ApJ...696....1D}) use one ordering;
the astrodust table uses another. Table~\ref{tab:jori} gives the mapping.

\begin{table}[!ht]
\centering
\caption{Orientation index conventions. The same three physical
orientations are numbered differently in two files from the same source.}
\label{tab:jori}
\begin{tabular}{lccc}
\toprule
Source & $k \perp a,\ E \parallel a$ & $k \perp a,\ E \perp a$ & $k \parallel a$ \\
\midrule
DF09 (\code{2009ApJ...696....1D}) & $j = 1$ & $j = 2$ & $j = 3$ \\
\code{q\_DH21Ad\_*}               & \code{jori} $= 2$ & \code{jori} $= 3$ & \code{jori} $= 1$ \\
\bottomrule
\end{tabular}
\end{table}

A third convention exists in the same tree: in
\code{qlib\_gra\_D16MGemt\_1.400}, \code{jori = 1} is the
\emph{random-orientation average}, not $k \parallel a$. This is
documented in the code at
\code{astrodust/sed/src/q\_graphite\_d16.f90:23--26}. Any reader routine
must therefore dispatch on the file identity, never on the index alone.

\subsection{HD23 alignment assumptions}

For reference, the alignment model adopted by HD23 is: oblate spheroids
with $b/a = 1.4$, porosity $0.20$, and
\begin{equation}
  \falign(a) \;=\; \frac{f_{\max}}{1 + (\aalign/a)^{\alpha}},
  \qquad
  \aalign = 0.0749\um,\quad \alpha = 1.80,\quad f_{\max} = 1.0,
  \label{eq:falign}
\end{equation}
with zero polarized contribution from PAHs. The axial ratio $b/a = 1.4$
is not a free aesthetic choice: it is the minimum departure from
sphericity that the observed starlight polarization efficiency integral
of DH21b will tolerate.

%====================================================================
\section{What alignment theory hands to radiative transfer}
\label{sec:alignment}

\subsection{The interface is one number}

Alignment theory, however elaborate, delivers exactly one scalar per
grain radius to the transfer solver: $R(a)$. It factorizes into internal
alignment (the angle $\beta$ between the short axis $\hat a_1$ and the
angular momentum $J$) and external alignment (the angle $\xi$ between $J$
and $B$),
\begin{equation}
  R \;=\; \left\langle
    \frac{3\cos^2\beta - 1}{2}\cdot\frac{3\cos^2\xi - 1}{2}
  \right\rangle,
  \label{eq:Rprod}
\end{equation}
following Hoang et al.\ (2024, ApJ 965, 183; \code{2024ApJ...965..183H},
Eq.~5). The practical reduced form used in that same paper (their Eq.~12)
is a two-state model,
\begin{equation}
  R = \begin{cases}
    0, & a < \aalign, \\[2pt]
    \fhighJ\,Q_X^{\rm high\text{-}J}
      + (1-\fhighJ)\,Q_X^{\rm low\text{-}J}, & a \geq \aalign,
  \end{cases}
  \qquad
  Q_X^{\rm high\text{-}J} = 1,\quad Q_X^{\rm low\text{-}J} \simeq 0.3.
  \label{eq:Rreduced}
\end{equation}

\subsection{The analytic alignment radius}

The reference analytic chain for \aalign\ is Hoang, Tram \& Lee (2021,
ApJ 908, 218; \code{2021ApJ...908..218H}). The RAT efficiency is
approximated by
\begin{equation}
  \bar Q_\Gamma \simeq
  \begin{cases}
    8\,(\bar\lambda/a)^{-3}, & a \lesssim \bar\lambda/2.7, \\
    0.4, & a \gtrsim \bar\lambda/2.7,
  \end{cases}
  \label{eq:qgamma}
\end{equation}
the infrared rotational damping coefficient by
\begin{equation}
  F_{\rm IR} \simeq 0.038\,U^{2/3}\,a_{-5}^{-1}\,n_3^{-1}\,T_1^{-1/2},
  \label{eq:fir}
\end{equation}
and the suprathermal rotation condition $\omega_{\rm RAT} = 3\,\omega_T$
then yields
\begin{equation}
  \aalign \simeq 0.055\,\hat\rho^{-1/7}
  \left(\frac{n_3 T_1}{\gamma_{-1} U}\right)^{2/7}
  \left(\frac{\bar\lambda}{1.2\um}\right)^{4/7}
  (1 + F_{\rm IR})^{2/7}\ \um .
  \label{eq:aalign}
\end{equation}
Inside a cloud the local field is reddened and attenuated, which is
captured by
\begin{equation}
  U = \frac{U_0}{1 + 0.42\,\Av^{1.22}},
  \qquad
  \bar\lambda = \bar\lambda_0\left(1 + 0.27\,\Av^{0.76}\right),
  \qquad
  \bar\lambda_0 = 1.3\um .
  \label{eq:avdep}
\end{equation}

\subsection{Rotational disruption}

Rotational disruption (RATD) removes large grains whose spin stress
exceeds the tensile strength $S_{\max}$. The important structural point
is that disruption defines a \emph{band}, not a cutoff: there is an upper
bound $a_{\rm disr,max}$ as well as \adisr, because sufficiently large
grains again spin too slowly to be disrupted. Implementing RATD as a
simple truncation of the size distribution is wrong.

\subsection{Magnetically enhanced alignment and its 2025 revision}
\label{sec:mrat}

Magnetically enhanced RAT alignment (MRAT) is parameterized by the
magnetic relaxation ratio $\delta_{\rm mag}$. Giang et al.\ (2023) adopt
a step function, $\fhighJ = 0.25$, $0.5$, $1$ for $\delta_m < 1$,
$1$--$10$, and $>10$ respectively. That last branch is the one to be
careful with.

Hoang (2025, ApJ 994, 115) revises the criterion to
\begin{equation}
  \delta_{\rm mag,cri} \;=\; (1 + F_{\rm IR})\,
  \frac{2 - q^{\max}}{q^{\max}},
  \label{eq:deltacri}
\end{equation}
showing that the familiar ``$\delta_m > 10$ implies perfect alignment''
rule was the special case $F_{\rm IR} \ll 1$. For $U \sim 10^3$ one has
$F_{\rm IR} \sim 60$, so the required $\delta_{\rm mag}$ rises by roughly
the same factor. Any implementation that keeps the 2016-era step function
will therefore overestimate alignment in strong radiation fields, exactly
where the polarized emission is brightest. The same paper finds that the
attainable maximum of \fhighJ\ is about $0.45$ for superparamagnetic
grains and about $0.22$ for ordinary paramagnetic grains, so the value
$1.0$ used by POLARIS is optimistic. It also proposes the environment
classifier $U/(n_1 T_2)$: values $\leq 1$ are collision-dominated, values
$>1$ are radiation-dominated, in which case the slow-rotation branch has
$f^{\rm slow}_{\rm high\text{-}J} = 0$.

Note the division of labor between these references: the authoritative
closed-form expressions for \aalign\ and \adisr\ remain those of Hoang et
al.\ (2021). Hoang (2025) does not supply new closed forms; it recasts
the problem in a dimensionless parameter space.

\subsection{Angle convention}

In this literature $\gamma$ denotes the angle between the \emph{mean
magnetic field and the line of sight}, not a plane-of-sky angle.
Polarization is maximal at $\gamma = 90^\circ$. Mixing this up with a
position angle is a common and quiet error.

\subsection{Which axis do grains align with?}
\label{sec:krat}

Everything above assumes the alignment axis is $\mathbf{B}$. That holds only
while Larmor precession is the fastest precession the grain undergoes.
Lazarian \& Hoang (2019, \code{2019ApJ...883..122L}) set the comparison out
explicitly. A grain with a Barnett magnetic moment precesses about
$\mathbf{B}$ on a timescale (their Eq.~4)
%
\begin{equation}
t_{\rm L} \;=\; \frac{4\pi}{5}\frac{g_e\mu_B}{\hbar}\,
   \rho_s\,s^{-2/3}\,a_{\rm eff}^{2}\,B^{-1}\chi(0)^{-1}
   \;\approx\; 1.3\;\hat\rho\;\hat s^{-2/3}\,a_{-5}^{2}\,
   \hat B^{-1}\hat\chi^{-1}\ \ {\rm yr},
\label{eq:tlarmor}
\end{equation}
%
with $\hat B = B/5\,\mu$G and $\hat\chi = \chi(0)/10^{-4}$. Anisotropic
radiation drives a competing precession about $\mathbf{k}$ (their
Eqs.~12--14),
%
\begin{equation}
t_{\rm rad,p} = \frac{2\pi}{\Omega_p},\qquad
\Omega_p = \frac{\bar u_{\rm rad}\,\bar\lambda\,a_{\rm eff}^{2}}
                {I_\parallel\,\omega}\,\gamma\,|\bar Q_\Gamma|,
\label{eq:tradp}
\end{equation}
%
\begin{equation}
t_{\rm rad,p} \approx 1.1\times10^{2}\;
  \hat\rho^{1/2}\hat s^{-1/3} a_{-5}^{1/2}\hat T_d^{1/2}
  \left(\frac{u_{\rm rad}}{u_{\rm ISRF}}\right)^{-1}
  \left(\frac{\bar\lambda}{1.2\,\mu\mathrm{m}}\right)^{-1}
  \left(\frac{\gamma|\bar Q_\Gamma|}{0.01}\right)^{-1}\ \ {\rm yr}.
\label{eq:tradp2}
\end{equation}
%
The criterion is simply which is shorter: if $t_{\rm rad,p} < t_{\rm L}$,
``the radiation direction rather than that of magnetic field becomes the axis
of alignment'' and the grain is in the k-RAT regime rather than B-RAT.

Two consequences matter for the present purpose. First, the paper's own
estimate is that both the radiative and the mechanical precession rates ``are
significantly smaller than the rate of Larmor precession,'' which ``ensures
that over most of the volume of the interstellar medium the alignment is
happening in respect to magnetic field.'' For a galaxy-scale model such as
NGC~891 the B-RAT assumption is therefore safe, and Eq.~(\ref{eq:tlarmor})
versus Eq.~(\ref{eq:tradp2}) can be evaluated as a diagnostic rather than
carried as a branch in the solver. The regime to watch is the immediate
vicinity of bright sources, where $u_{\rm rad}/u_{\rm ISRF}$ is large enough
to push $t_{\rm rad,p}$ below $t_{\rm L}$.

Second, $t_{\rm L} \propto \chi(0)^{-1}$, so iron inclusions that raise the
susceptibility also push the transition further away: the same magnetic
enhancement that produces efficient MRAT alignment
(\S\ref{sec:mrat}) simultaneously protects the B-RAT assumption. Conversely,
detecting a k-RAT signature would constrain $\chi(0)$, which is the
observational program that paper proposes.

Note also that Eq.~(\ref{eq:tradp2}) uses
$|\bar Q_\Gamma| \approx 2.3(\lambda/a_{\rm eff})^{-3}$ for
$\lambda > 1.8 a_{\rm eff}$ and $0.4$ otherwise, whose coefficient and
transition point differ from the parameterization quoted in
Eq.~(\ref{eq:qgamma}). The two are alternative fits to the same DDA results
and agree to within tens of percent; use one consistently rather than mixing
them.

\subsection{The polarization hole}

The mechanism is a chain of three steps, each of which the formulas above
already contain: increasing \Av\ lowers $U$ (Eq.~\ref{eq:avdep}), lower
$U$ raises \aalign\ (Eq.~\ref{eq:aalign}), and once \aalign\ exceeds the
largest grain present, alignment vanishes and polarized emission goes to
zero. The complete-loss condition scales as
$A_{V,\max} \simeq 24.7\,(\cdots)$, evaluating to roughly $15.9$, $6.9$,
and $3.0$ mag for $n_{\rm H} = 10^4$, $10^5$, $10^6\ {\rm cm}^{-3}$ with
$a_{\max} = 0.25$\um. Grain growth pushes the transition to larger \Av,
which is why the observed depth of the hole is a joint constraint on
alignment and on growth rather than on either alone.

%====================================================================
\section{Polarized transfer and its implementation}
\label{sec:rt}

\subsection{The Stokes transfer equation}

The transfer of the Stokes vector through an aligned medium is governed
by a $4\times4$ extinction matrix. In the form given by POLARIS~I
(Reissl, Wolf \& Brauer 2016, A\&A 593, A87;
\code{2016A\&A...593A..87R}, their Eq.~7),
\begin{equation}
  \frac{d}{ds}
  \begin{pmatrix} I \\ Q \\ U \\ V \end{pmatrix}
  = -\,n_d
  \begin{pmatrix}
    \Cext & \Cpol\cos2\varphi & \Cpol\sin2\varphi & 0 \\
    \Cpol\cos2\varphi & \Cext & 0 & \Ccirc\sin2\varphi \\
    \Cpol\sin2\varphi & 0 & \Cext & -\Ccirc\cos2\varphi \\
    0 & -\Ccirc\sin2\varphi & \Ccirc\cos2\varphi & \Cext
  \end{pmatrix}
  \begin{pmatrix} I \\ Q \\ U \\ V \end{pmatrix}
  + j,
  \label{eq:transfer}
\end{equation}
where $\varphi$ is the position angle of the projected field and
$j$ is the polarized emissivity.

\paragraph{Typographical error in the published matrix.} The printed
$(4,4)$ element in that paper is $\Cpol$; the correct entry is $\Cext$,
as written in Eq.~(\ref{eq:transfer}) above. The reasoning is structural:
$V$ must suffer the same total extinction as $I$, $Q$, $U$, and only
$\Ccirc$, appearing antisymmetrically, may rotate $U$ into $V$. The
POLARIS source code implements it correctly, so this is an error in the
paper rather than in the code.

\subsection{How POLARIS assembles the matrix}

The assembly happens in \code{DustMixture.cpp::calcEmissivityExt}, whose
non-zero off-diagonal assignments are:

\begin{lstlisting}[style=cpp]
K[0][1] = K[1][0] =  Cpol  * cos2phi;
K[0][2] = K[2][0] =  Cpol  * sin2phi;
K[1][3] =            Ccirc * sin2phi;
K[2][3] =           -Ccirc * cos2phi;
K[3][1] =           -Ccirc * sin2phi;
K[3][2] =            Ccirc * cos2phi;
\end{lstlisting}

The cross sections themselves absorb the alignment factor in
\code{DustComponent.cpp::calcCrossSections}, as
\begin{align}
  \Cpol  &= \tfrac12\,R\,(\Cext^{\perp} - \Cext^{\parallel})\,\sin^2\theta,
    \label{eq:cpolR} \\
  C_{\rm pabs} &= \tfrac12\,R\,(\Cabs^{\perp} - \Cabs^{\parallel})\,\sin^2\theta.
    \label{eq:cpabsR}
\end{align}
A single scalar $R$ thus scales the dichroic and the birefringent
amplitudes simultaneously, which is the concrete expression of the
statement in \S\ref{sec:alignment} that the alignment interface is one
number.

\subsection{Prefer the closed-form solution to adaptive integration}

If the density and the alignment direction are constant within a cell,
$\exp(-Kns)$ has a closed form (Peest et al.\ 2023,
\code{2023A\&A...673A.112P}): $I$ and $Q$ evolve through
$\cosh$ and $\sinh$ of $\Cpol\,n s$, $U$ and $V$ through $\cos$ and
$\sin$ of $\Ccirc\,n s$, with an overall factor $e^{-\Cext n s}$. POLARIS
instead integrates Eq.~(\ref{eq:transfer}) with an adaptive
Runge--Kutta--Fehlberg 4(5) scheme, applying the error test to each of
the four Stokes components with $\epsilon_{\rm err} = 10^{-8}$. For a
cell-based grid the closed form is strictly better: it is exact, it is
cheaper, and it removes a tolerance parameter. Peest et al.\ (2023) also
publish analytic benchmark problems, which their MCpol implementation
reproduces to within $0.1\%$; those benchmarks should be adopted as
regression tests.

\subsection{The three-pass structure and what it buys}

POLARIS runs three passes, and the division of labor is the single most
useful architectural lesson available:

\begin{enumerate}
\item Monte Carlo thermal equilibrium, with alignment ignored.
\item Monte Carlo radiation field, storing a \emph{vector}
      $u_\lambda$ and the anisotropy measure
      $\gamma_\lambda = |\sum u_i| / \sum|u_i|$.
\item Observables: alignment switched on, Stokes transfer along rays.
\end{enumerate}

During the Monte Carlo passes the alignment is deliberately randomized
(\code{DustMixture.cpp:652--656}). The consequence is worth stating
explicitly, because it removes what looks like the largest obstacle:
\emph{polarized observables do not require the Monte Carlo transport
itself to be polarized}. What is given up is the feedback of dichroic
extinction on the radiation field, which is an approximation but a
controlled one.

\subsection{The radiation field must be vectorial}
\label{sec:vecfield}

RAT torque depends on the anisotropy of the local radiation field, so a
scalar mean intensity is insufficient: alignment simply cannot be
computed from $J_\lambda$ alone. What is required per cell and
wavelength is the vector sum as well as the scalar one,
%
\begin{equation}
\bar u_\lambda = \sum_i |\vec u_{\lambda,i}|,\qquad
\gamma_\lambda = \frac{\left|\sum_i \vec u_{\lambda,i}\right|}
                        {\sum_i |\vec u_{\lambda,i}|} \in [0,1],
\label{eq:gammalam}
\end{equation}
%
with $\gamma_\lambda$ the anisotropy degree that enters the torque. This has a hard corollary. The modified
random walk (MRW) destroys directional information by construction and
therefore cannot be used in the pass that feeds alignment. It remains
available in the thermal pass.

\subsection{Torque efficiency tables are not needed}

A pleasant surprise: POLARIS's \code{calcAlignedRadii} has the tabulated
route (\code{getQrat}) commented out and uses the power-law approximation
of Eq.~(\ref{eq:qgamma}) instead. Equivalent results are therefore
obtainable without computing or storing a T-matrix torque-efficiency
table, which removes what would otherwise be a substantial Phase-A style
computation.

\subsection{Stochastic heating and polarization}
\label{sec:stoch}

POLARIS supports the combination, using the Camps et al.\ (2015)
formulation, and the structure is simple: the temperature probability
weighting $P(T)$ multiplies $I$ and $Q,U$ identically. Two consequences
follow. Within a single grain size the polarization fraction is
\emph{independent} of $P(T)$; the temperature distribution affects the net
polarization only through the relative weighting of different sizes.

Physically the two populations barely overlap in any case. Alignment sets
in near $\aalign \simeq 0.075$\um, whereas stochastic heating dominates
for $a \lesssim 0.01$--$0.05$\um. The observational counterpart is the
non-detection of polarization in the 3.4\um\ C--H feature, $0.06\% \pm
0.13\%$ (Chiar et al.\ 2006, ApJ 651, 268; \code{2006ApJ...651..268C}),
consistent with the HD23 assumption of zero polarized PAH emission.

Nevertheless the correct treatment is essentially free here, and that is
an argument specific to this code base. SEDust already solves $P(T)$;
carrying the polarized emission through the same integral costs one
additional accumulator. It is worth noting that the Hoang group papers
all use equilibrium temperatures only; Tram et al.\ (2025, GRADE-POL)
states explicitly that the temperature distribution was discarded. Doing
it properly would be a genuine, if modest, improvement over the reference
implementations.

%====================================================================
\section{Current state of the codes}
\label{sec:codes}

\subsection{SEDust}

\emph{Status note (2026 July 19): stage~1 of Table~\ref{tab:stages} is
implemented and verified. The paragraphs below are kept in the past tense
where they describe the starting point, and the achieved state is
summarized at the end of this subsection. \S\ref{sec:worklist} likewise
records what was actually done.}

SEDust was in better shape than expected. The polarized cross-section
table \code{q\_DH21Ad} was already present with all three orientation
blocks (\S\ref{sec:optics}), and its correctness has been verified
against the HD23 release. The \code{tmatrix/} directory links only
\code{tmd\_one.f}, the random-orientation average driver; the
fixed-orientation code \code{ampld.lp.f} is present but unused. As long
as the astrodust model is the target, no T-matrix run is needed for the
cross sections; a run is needed only for the scattering matrix
(\S\ref{sec:optical}, now done for five optical bands), for exploring
other axial ratios or compositions, or for the fixed-orientation
amplitude matrix that circular polarization would require (\S\ref{sec:optics}): the released table
carries $Q_{\rm ext}$, $Q_{\rm abs}$ and $Q_{\rm sca}$ only, so \Ccirc\
cannot be built from it.

The structural entry point was \code{grain\_pop\_t} in
\code{sed/src/dust\_model\_mod.f90}: adding a \code{Cpol(NLAM,NA)}
component there propagates naturally through the existing population
loop. On the emission side, one polarized accumulator alongside the
existing one in \code{sed\_astrodust.f90} completed the change.

\paragraph{What SEDust now provides.} The orientation-resolved reader is
\code{sed/src/q\_table\_jori.f90}. \code{sed\_init} builds two polarized
cross sections on the size-distribution grid, \code{Cpol} (polarized
absorption, which drives the polarized emission) and \code{Cpol\_ext}
(polarized extinction, which drives the dichroic term of the extinction
matrix), together with $\falign(a)$ from the HD23 fit. Three exports
follow:
\begin{itemize}
\item \code{dust\_emission} takes an optional \code{lamI\_pol} of the
      same shape and units as \code{lamI\_total}, returning the
      \emph{intrinsic} polarized emission;
\item \code{../data/kext\_astrodust\_MW.dat} carries the polarized
      extinction per H as an eighth column, appended so that readers of
      the first seven columns are unaffected;
\item \code{calc\_polext.x} regenerates the comparison against the
      released \code{polarized\_extinction.dat}.
\end{itemize}
The reproduction is a median $|{\rm ours/ref}-1|$ of $0.0296\%$ over the
992 points with $\lambda \ge 0.11\um$, and the intrinsic polarization
fraction of the emission is $19.15\%$ at $850\um$ and $17.15\%$ at
$154\um$. The calling interface is specified in the SEDust user manual;
\code{docs/mocafe\_polarization\_handoff.md} records the decisions taken
while building this and the limits of the released data.

\subsection{MoCafe v2.00}

MoCafe already carries the polarization state. \code{photon\_type} has
\code{I,Q,U,V} (\code{define.f90:117}), there is a \code{use\_stokes}
switch (line 213), and \code{scatter\_dust\_stokes} interpolates
$S_{11}$, $S_{12}$, $S_{33}$, $S_{34}$. SEDust is embedded through
\code{dustemis\_mod}, and the code passes the Camps et al.\ (2015) SHG
benchmark. The task is already on record: the roadmap contains
\code{[TODO, L] Polarized thermal emission from aligned grains.
Long-term.}

Four things are missing:
\begin{enumerate}
\item No magnetic field exists anywhere (a source-wide grep returns zero
      hits). This is the largest single gap.
\item Extinction is scalar: \code{tau = rhokap * length}.
\item \code{jtally\_mod} accumulates a \emph{scalar} $J_\lambda$ only,
      and must be vectorized before any self-consistent alignment
      calculation is possible.
\item There is no polarized emission source term.
\end{enumerate}

\subsection{Reference sources in the tree}

Two POLARIS checkouts are available for reference:
\code{Grain/POLARIS/} (branch \code{master-basic}, commit
\code{03280f5}) and \code{Grain/NewPOLARIS/} (branch \code{main}, commit
\code{e508e89}; tag \code{v1} = \code{a4fcb4d} corresponds to Zenodo
\code{10.5281/zenodo.14199014}, the published version, with \code{main}
ahead only in the README).

NewPOLARIS is the MRAT and RATD extension of Giang, Hoang \& Kim (2023).
It is \emph{not} merged upstream and was forked from an older layout
(\code{Dust.cpp}/\code{Dust.h}), so a direct diff against the baseline is
not possible; only the physics can be ported. What it adds: the
\code{ALIG\_MRAT}, \code{ALIG\_kRAT}, \code{ALIG\_INELASTIC} flags; the
keywords \code{<number\_cluster>}, \code{<volume\_filling\_cluster>},
\code{<iron\_fraction>}, \code{<change\_f\_highJ>}, \code{<inelascity>},
\code{<disrupt> RATD}; and grid fields holding critical radii for each
cell (\code{GRIDadisr}, \code{GRIDamaxJB\_Lar},
\code{GRIDadg\_lower/upper}), which moves the expensive physics into a
preprocessing step. The stock POLARIS has no MRAT, and its
\code{ALIG\_KRAT} is a dead flag: it is connected to neither the parser
nor any physics path.

%====================================================================
\section{Staged proposal}
\label{sec:stages}

\begin{table}[!ht]
\centering
\caption{Proposed stages, with the status as of 2026 July 19. Stage~2 is
the point at which quantitative comparison against SOFIA/HAWC+ becomes
possible; stage~1 no longer blocks it.}
\label{tab:stages}
\begin{tabularx}{\textwidth}{c l c c X}
\toprule
Stage & Scope & Effort & Status & Content \\
\midrule
1 & SEDust cross sections & days, low risk & done &
  \code{Cpol}, \code{Cpol\_ext} and \falign\ built in \code{sed\_init};
  polarized emission accumulator in \code{sed\_grain\_loop}; optional
  \code{lamI\_pol} on \code{dust\_emission}; polarized extinction
  exported as column~8 of \code{kext\_astrodust\_MW.dat}. Released
  \code{polarized\_extinction.dat} reproduced to a median of
  $0.0296\%$; intrinsic polarization fraction $19.15\%$ at $850\um$. \\
2 & MoCafe ray tracing & moderate & next &
  Polarized emission in ray tracing and peel-off, using the Peest
  closed form. Needs no scattering matrix: the far-infrared albedo is
  near zero. \\
3 & MoCafe forward MC & moderate & open &
  Dichroic extinction; keep scalar $\tau$ for path sampling and correct
  by weight. Column~8 supplies \Cpol. \\
4 & Self-consistent alignment & large & open &
  Vectorial radiation-field pass; MRW disabled; memory $\times 4$. \\
\bottomrule
\end{tabularx}
\end{table}

\paragraph{Stage 1 (done).} \code{Cpol}, \code{Cpol\_ext} and \falign\
are built in SEDust, with a single polarized emission accumulator
alongside the existing one. Nothing in the solver core changed, and the
three astrodust SEDs plus the DL07 SED regress byte-identically. The
first validation step, reproducing the released
\code{polarized\_extinction.dat}, passes at a median of $0.0296\%$. The
remaining checks, $p_{\rm submm}$ against Planck, $R_{S/V}$, and the
Serkowski parameters $\lambda_{\max}$ and $K$, are now possible from
column~8 of \code{kext\_astrodust\_MW.dat} and remain to be done.

\paragraph{Stage 2.} Polarized emission in the MoCafe ray tracing and
peel-off paths, integrated with the closed-form solution of
\S\ref{sec:rt}. This is the milestone worth aiming at first: it already
allows a quantitative comparison with the HAWC+ 154\um\ maps, which did
not exist in 2018.

\paragraph{Stage 3.} Dichroic extinction in the forward Monte Carlo. The
recommended approach keeps the scalar $\tau$ for path-length sampling and
applies a weight correction, which sidesteps the circularity analyzed by
Baes et al.\ (2019). At this stage the empirical \code{dichroic.f90}
formula is replaced by Eq.~(\ref{eq:tanh}) with a wavelength-dependent
$\Cpol$.

\paragraph{Stage 4.} Alignment computed internally. This requires the
vectorial radiation-field pass, forbids MRW in that pass, and raises the
radiation-field memory by a factor of four. Until then the sensible
choice is the SKIRT approach: treat $\falign(a)$ as an input function,
with the HD23 fit of Eq.~(\ref{eq:falign}) as the default.

%====================================================================
\section{What the update actually touches}
\label{sec:worklist}

The stages above are stated in physical terms. This section translates
them into the specific data structures and routines that would change,
based on the state of the two trees as of 2026 July. Line numbers are
given where they pin down a unique insertion point; they will drift.

\subsection{SEDust, stage 1 (implemented)}

\emph{This subsection was written as a plan and is kept for the reasoning
it records. What was actually built follows each paragraph in italics.
The resulting calling interface is specified in the SEDust user manual.}

\paragraph{Read the polarized cross sections.} The existing loader
\code{load\_q\_table} (\code{sed/src/q\_table.f90}) reads the flat table
produced by \code{run\_tmatrix}, whose columns are
\code{lambda a\_eff Qext Qabs Qsca albedo g flag}. The HD21 file has a
different shape: three orientation blocks, each
$1129 \times 169$, written as
\code{((Q(jw,jr,jori),jw=0,1128),jr=0,168),jori=1,3} for each of
$Q_{\rm ext}$, $Q_{\rm abs}$, $Q_{\rm sca}$ in turn. Two options exist:
add a second reader for the HD21 layout, or extend the existing table
format with a \code{jori} column and regenerate. The first is faster to
get running; the second keeps one reader for both the shipped table and
any future T-matrix computation, and is the better end state. Whichever
is chosen, form
%
\begin{equation}
Q_{\rm pol}^{\rm abs}(\lambda,a)
  = \tfrac{1}{2}\left[Q_{\rm abs}^{(jori=3)} - Q_{\rm abs}^{(jori=2)}\right],
\qquad
Q_{\rm pol}^{\rm ext} \ \hbox{likewise from } Q_{\rm ext},
\label{eq:qpolfromjori}
\end{equation}
%
and watch the index convention of Table~\ref{tab:jori}.

\emph{Done as the first option: \code{sed/src/q\_table\_jori.f90} is a
second reader for the HD21 layout, leaving \code{q\_table.f90}
untouched. It decompresses the shipped \code{.gz} through \code{gzip
-dc} into a scratch copy, and exposes \code{qpol\_ext}, \code{qpol\_abs},
the random-orientation combinations \code{qran\_*}, the grids
\code{lam\_j}, \code{aeff\_j}, and \code{falign\_hd23(a)}.}

\paragraph{Carry it on the population.} \code{grain\_pop\_t}
(\code{sed/src/dust\_model\_mod.f90}) currently holds \code{aeff},
\code{dn}, \code{Cabs}, \code{Csca}, \code{kappB}, \code{H} and
\code{kappCMB}. Add
%
\begin{lstlisting}[language=Fortran]
real(wp), allocatable :: Cpol(:,:)   ! (NLAM, NA)  polarized absorption
real(wp), allocatable :: falign(:)   ! (NA)        alignment efficiency
\end{lstlisting}
%
and free them in \code{free\_pop} alongside the rest. Filling
\code{falign} from Eq.~(\ref{eq:falign}) needs only \code{aeff}, which
the population already carries after the earlier fix that added it.

\emph{Done, with \code{set\_pop} taking \code{Cpol} and \code{falign} as
optional arguments so that populations without polarized optics simply
leave them unallocated. The module-level arrays behind them are
\code{Cpol} and \code{Cpol\_ext}; the short name stays on the absorption
one because the emission path is written around it.}

\paragraph{Accumulate the polarized emission.} The emission loop in
\code{sed\_grain\_loop} (\code{sed/src/sed\_astrodust.f90}) has four
accumulation sites, at roughly lines 757, 762, 926 and 938, of the form
%
\begin{lstlisting}[language=Fortran]
Jout_local = Jout_local + dn_pop(ir) * P(ii) * Cabs_pop(:,ir) * spec
\end{lstlisting}
%
Each gains a companion line that differs only in the cross section and
the alignment weight:
%
\begin{lstlisting}[language=Fortran]
Jpol_local = Jpol_local + dn_pop(ir) * P(ii) * Cpol_pop(:,ir) &
                                     * falign_pop(ir) * spec
\end{lstlisting}
%
This is the point made in \S\ref{sec:stoch}: because the temperature
weight multiplies both accumulators identically, the exact $P(T)$
treatment comes for free, and there is no reason to approximate the
polarized emission with an equilibrium temperature. The two sites at 757
and 938 are the equilibrium branches and need the same companion.

\emph{Done at every accumulation site, guarded by a single \code{do\_pol}
flag that is true only when the caller asked for the polarized output and
the population carries both arrays.}

\paragraph{Widen the public interface.} \code{dust\_emission}
(\code{sed/src/sed\_astrodust.f90:2143}) currently returns
\code{lamI\_total} and optionally \code{lamI\_chan}. Add an optional
\code{lamI\_pol} of the same shape as \code{lamI\_total}. Making it
optional keeps every existing caller valid, which matters because
MoCafe links this as a library; the same discipline was used for the
\code{status} argument. Re-export through \code{dust\_lib}.

\emph{Done. The signature is now}
\code{dust\_emission(m, J\_lam, lamI\_total [, lamI\_chan] [, status]
[, lamI\_pol])}\emph{. The polarized \emph{extinction} took a different
route: rather than a second library accessor, it is written as column~8
of \code{../data/kext\_astrodust\_MW.dat}, because that file is already
how a host consumes extinction and \code{dust\_lib} exposes no
extinction surface for an accessor to sit beside.}

Note that what leaves SEDust is the \emph{intrinsic} polarized
emissivity, that is, the quantity before the geometric factor
$\sin^2\gamma$ and before any turbulent reduction. Those depend on the
local field direction and belong to the RT side. Keeping the split here
is what allows a single pair of tables to serve all viewing geometries.

\paragraph{Verify before going further.} Reproduce
\code{data/release/polarized\_extinction.dat} by combining
Eq.~(\ref{eq:cpol}) with the alignment fit of Eq.~(\ref{eq:falign})
and integrating over the size distribution. This exercises the orientation indexing, the sign
convention, the $\falign$ fit and the size integration in one shot. It
has already been shown to work at 0.55\um; extending the check across
the full wavelength grid is the natural first commit. Only then are the
$p_{\rm submm}$, $R_{S/V}$ and Serkowski comparisons of
\S\ref{sec:validation} meaningful.

\emph{Done by \code{calc\_polext.x}, which reproduces the release to a
median of $0.0296\%$ and a maximum of $0.94\%$ over the 992 points with
$\lambda \ge 0.11\um$ (the reference itself turns negative below that).
Column~8 of the extinction table agrees with that program to the seven
digits the ASCII format carries, so the two routes to the same quantity
are consistent. The $p_{\rm submm}$, $R_{S/V}$ and Serkowski comparisons
are still open.}

\subsection{MoCafe, stage 2}

\paragraph{A magnetic field has to exist.} No field of any kind appears
in the current source. A minimal implementation adds a field direction
for each cell, either from an analytic model or from an MHD snapshot,
alongside the existing density. For NGC~891 the analytic route is the
starting point, but the observational result of Kim et al.\ (2023) means
a turbulent component is not optional: a smooth field cannot produce the
low polarization fraction and narrow dispersion they report. In practice
this means either resolving the turbulence on the grid or carrying a
reduction factor $F_{\rm turb}$ per cell, as GRADE-POL does.

\paragraph{Emissivity becomes a Stokes vector.}
\code{compute\_dustemis} (\code{src/dustemis\_mod.f90}) currently
obtains a scalar spectrum from \code{dust\_emission} or
\code{dust\_emission\_interp} and stores one emissivity array per cell.
With \code{lamI\_pol} available it stores, for each cell,
%
\begin{equation}
j_I = \hbox{(as now)},\qquad
\binom{j_Q}{j_U} = \lamIpol \, \sin^2\gamma \, F_{\rm turb}
   \binom{\cos 2\varphi}{\sin 2\varphi},
\label{eq:jstokes}
\end{equation}
%
where $\gamma$ is the angle between the local field and the line of
sight to the observer and $\varphi$ is the position angle of the
field projected on the sky. Both follow from the field direction and the
observer direction, so no new physics enters here.

There is a structural consequence worth noting: $\gamma$ and $\varphi$
depend on the viewing direction, so $j_Q$ and $j_U$ are not properties
of the cell alone in the way $j_I$ is. For a single observer this is
immaterial. For the all-sky output of \code{allsky\_mod} it means the
angles must be evaluated per direction rather than cached per cell.

\paragraph{Integrate the Stokes vector along the ray.} The ray-tracing
and peel-off paths (\code{src/peelingoff\_mod.f90},
\code{src/allsky\_mod.f90}) currently accumulate a scalar optical depth
through \code{raytrace\_tau} and apply $e^{-\tau}$. Replace this, for
the polarized output only, by the segment-by-segment application of the
closed-form solution of \S\ref{sec:rt}: within a segment of constant
density and field direction the Stokes vector transforms by
$\cosh$ and $\sinh$ of $\Cpol n s$ in $(I,Q)$ and by $\cos$ and $\sin$
of $C_{\rm circ} n s$ in $(U,V)$, all scaled by $e^{-\Cext n s}$, with a
rotation into and out of the local field frame at each segment boundary.
The existing segment walk already provides the geometry; what changes is
what is done with each segment.

At this point the model produces polarized emission maps that can be
compared with the HAWC+ 154\um\ data directly. Nothing in the forward
Monte Carlo has been touched.

\subsection{MoCafe, stage 3}

The forward Monte Carlo is where dichroism becomes awkward, because the
optical depth that governs path sampling is no longer a scalar. The
recommendation of \S\ref{sec:stages} is to keep sampling path lengths
from the scalar $\Cext$ and to apply the polarization as a weight
correction on the surviving packet, which preserves the sampling
distribution and avoids the circularity that Baes et al.\ (2019)
analyze. The physical cost is that dichroic extinction does not feed
back on the radiation field; whether that matters can be measured by
comparing stage-3 output with stage-2 output, which is a reason to build
them in this order.

This is also the point at which \code{Galaxy\_pol/dichroic.f90} would be
retired. Its replacement is Eq.~(\ref{eq:tanh}) with $\Cpol$ from
Eq.~(\ref{eq:qpolfromjori}), which restores the wavelength dependence
that the empirical form cannot represent.

\subsection{MoCafe, stage 4}

Only here does the radiation-field tally change.
\code{jtally\_mod} stores \code{jt\_sum(nlambda,ncell)}, a scalar mean
intensity. Radiative torque alignment needs the direction of the
radiation as well, so the tally becomes four numbers per cell and
wavelength: the scalar sum and the three components of
$\sum_i \vec{u}_{\lambda,i}$, from which the anisotropy degree
$\gamma_\lambda$ of Eq.~(\ref{eq:gammalam}) follows. The memory note
already in the roadmap, that \code{jt\_sum} is held privately per rank,
becomes four times more pressing.

Two further constraints come with it. The modified random walk cannot be
used in this pass, since it destroys the directional information. And
the alignment solve itself is a preprocessing step over cells, not an
inner-loop cost: as NewPOLARIS shows, the sensible structure is to
compute $a_{\rm align}$, $a_{\rm disr}$ and the magnetic radii once per
cell, store them, and have the cross-section evaluation do nothing but
compare grain size against stored radii.

\subsection{Order of work and what each stage buys}

Stage 1 is self-contained inside SEDust and verifiable against a
published file. Stage 2 produces the first result that could not have
been obtained in 2018, because the HAWC+ maps did not exist. Stage 3
removes the empirical formula that is the weakest point of the 2018
model. Stage 4 replaces the assumed $\falign(a)$ with computed
alignment, and is the only stage that changes the Monte Carlo itself.

A useful property of this ordering is that each stage is independently
publishable and independently falsifiable: stage 1 against
\code{polarized\_extinction.dat}, stage 2 against HAWC+, stage 3 against
stage 2, and stage 4 against the fixed-$\falign$ result of stage 2.

%====================================================================
\section{What optical polarized RT additionally requires}
\label{sec:optical}

Everything up to this point is written with the far infrared in view,
where the albedo is near zero and a scattering matrix can be skipped
(\S\ref{sec:stages}, stage~2). Optical starlight polarization is the
case the 2018 model actually addressed, and it needs strictly more.
This section collects the additional requirements in one place, because
they are currently scattered: the fixed-orientation amplitude code is
mentioned in \S\ref{sec:codes}, \Ccirc\ in \S\ref{sec:optics} and
\S\ref{sec:rt}, and the MPFA in \S\ref{sec:optics}.

The immediate goal is to update Seon (2018), which worked at a single
wavelength ($V$ band, \S\ref{sec:motivation}). Concretely, the target is
to replace the empirical $P_V = 2.7\,\tau_V^{0.48}$ of
Eq.~(\ref{eq:whittet}) by dichroism computed from cross sections. A
plausible extension is one or two further optical bands, which is a
different thing from a full wavelength grid: the far infrared is not in
the present plan at all, and neither is the ultraviolet. The
requirements below are ordered on that basis, and items that only bind
once the ultraviolet or the far infrared enters are marked as such.
Several requirements that look formidable when stated for a full
wavelength grid stay modest across two or three optical bands, and the
ordering exploits that.

Two consequences of the wider optical scope are worth stating at the
outset. A scattering matrix would be needed at each band rather than at
one wavelength, which multiplies the computation but not the method: a
handful of bands is still far from a grid. And once more than one band
is available the wavelength dependence of the polarization becomes an
observable in its own right, which is exactly what the empirical formula
of Eq.~(\ref{eq:whittet}) cannot reproduce, so the case for replacing it
grows stronger rather than weaker.

\subsection{The stage table presupposes data that does not exist yet}

Table~\ref{tab:stages} assigns stages~3 and~4 to MoCafe, which is
correct as far as transport is concerned but incomplete as a plan.
Those stages consume optical quantities that SEDust did not originally
produce, and in one case that no released file contains. Stage~1
delivered \Cpol\ and \code{Cpol\_ext}, sufficient for polarized emission
and for dichroic extinction; the random-orientation scattering matrix
has since been added for five optical bands. Table~\ref{tab:optical}
lists the supply side item by item, with the priority each carries for
the single-wavelength update.

\begin{table}[!ht]
\centering
\caption{Ingredients that far-infrared polarized transfer can do
without and optical polarized transfer cannot. ``Supplier'' is the code
that must produce the quantity; ``Priority'' is judged against the
single-wavelength ($V$ band) goal, not against a full wavelength grid.}
\label{tab:optical}
\begin{tabularx}{\textwidth}{l l l X}
\toprule
Ingredient & Supplier & Priority & Status, and why optical needs it \\
\midrule
$Z(\theta)$ at $V$, random orientation &
  SEDust & \textbf{done} &
  Computed by \code{run\_scatmat.x} and stored for five optical bands,
  \textit{UBVRI}. Scattering dominates the optical opacity, so the phase matrix
  sets the scattered polarization; in the far infrared the albedo is
  near zero and $Z$ never enters. \\
\falign\ functional form &
  either & high &
  HD23 fit assumed. Optical polarization is carried by grains near
  $a \sim \aalign$, so the shape of the cutoff is a first-order effect,
  and it is what sets how the polarization varies from band to band;
  far-infrared polarization comes from grains well above it, where all
  forms agree. \\
$Z(\lambda,\theta)$, aligned orientation &
  SEDust & later &
  \code{ampld.lp.f} present but unused. Aligned spheroids scatter
  anisotropically with respect to $\mathbf{B}$, which spherical Mie
  cannot represent. \\
\Ccirc &
  SEDust & later &
  Absent from the release. Birefringence converts $U$ into $V$; the
  2018 model had none either, so this matters only once a circular
  observable is in scope. \\
Averaging convention &
  SEDust & low &
  Mixed, exact average for \Cabs\ vs.\ MPFA for \Cpol. The disagreement
  is $0.080\%$ median across the optical, so it does not bind on a
  programme confined to optical bands. It becomes an obstacle in the
  ultraviolet, where the median is $0.21\%$ and the maximum $9.53\%$.
  Deliberately deferred in v1.20; see the paragraph below. \\
\bottomrule
\end{tabularx}
\end{table}

\paragraph{Scattering matrix, random orientation (available).} This item
is done. The T-matrix driver in the tree, \code{tmatrix/src/tmd\_one.f},
was originally derived from Mishchenko's \code{tmd.lp.f} by wrapping the
implicit main program as a subroutine, and in the process the final
\code{MATR} call was removed and the angular expansion coefficients were
dropped from the returned quantities. The \code{MATR} routine itself
remained in the file, and the call has since been restored through the
entry point \code{TMD\_ONE\_SCATMAT}.

\code{tmatrix/driver/run\_scatmat.f90} uses it to accumulate the
single-size matrices over the HD23 astrodust size distribution with
$C_{\rm sca}\,{\rm d}n/n_{\rm H}$ weighting and to expand the result onto
a scattering-angle grid. The output is
\code{tmatrix/output/scatmat\_astrodust\_P0.20\_Fe0.00\_1.400.dat}: one
block per wavelength, $181$ rows at $\theta = 0, 1, \ldots, 180^\circ$,
with the six independent elements a randomly oriented particle with a
plane of symmetry possesses, in the column order $F_{11}, F_{22},
F_{33}, F_{44}, F_{12}, F_{34}$. The normalization is Mishchenko's,
$\tfrac{1}{2}\int_{-1}^{+1} F_{11}\,{\rm d}\cos\theta = 1$, so absolute
units follow by multiplying by the $C_{\rm sca}$ quoted in each block
header. This already improves on the 2018 model, whose phase matrix was
spherical Mie.

Five wavelengths are stored, $0.36$, $0.44$, $0.55$, $0.64$ and
$0.79\um$, approximately \textit{UBVRI}. That covers the scope of the update as
planned --- the $V$ band first, with one or two more optical bands
possible, and no far-infrared or ultraviolet use foreseen. The full
1129-point sweep costs about 1.7 hours and would produce mostly
wavelengths outside that scope, so it was not run. It requires no code
change should the scope widen: \code{./run\_scatmat.x all} computes the
whole grid, or an explicit list of wavelengths a subset, at about
$5.5\,$s per wavelength.

What remains missing is the \emph{aligned} scattering matrix, and
\Ccirc; both are treated in the two paragraphs below. The distinction
matters: the random-orientation matrix answers how an unaligned grain
population scatters, which is what a first optical model needs for the
scattering source term, but it carries no information about the
orientation of the grains with respect to $\mathbf{B}$.

\paragraph{The cutoff shape is the physics specific to the optical
(needed).} Eq.~(\ref{eq:falign}) is a fitted function, and
\S\ref{sec:stages} treats it as an input rather than a computed
quantity. In the far infrared that is harmless, because the emission is
dominated by grains several times larger than \aalign, where any
reasonable form has saturated at $f_{\max}$. At $V$ the polarization is
produced by grains of size comparable to the wavelength, which is the
same size range as \aalign\ itself, so the result responds directly to
whether the cutoff is a step or a smooth roll-off. The observational
expression of this is that Serkowski's $\lambda_{\max}$ is set largely
by the alignment size threshold rather than by the size distribution.
Truong et al.\ (\code{2025ApJ...981...83T}) make the same point from the
other side, noting that the step and smooth forms are distinguishable
only in the optical. This item does not become easier at a single
wavelength and keeps its priority.

\paragraph{Scattering matrix, aligned orientation (later).} The full
orientation-dependent amplitude matrix requires \code{ampld.lp.f}, which
is present in \code{tmatrix/src/} but excluded from the build. Stated
for a wavelength grid the cost is storage rather than algorithm, since
the table acquires a wavelength, a scattering angle, an orientation and
$16$ matrix elements. At a single wavelength the wavelength axis drops
out and the table is far smaller, which is worth remembering before this
item is dismissed as infeasible. There is no shortcut through published
data: the HD21 release carries integrated $Q_{\rm ext}$, $Q_{\rm abs}$
and $Q_{\rm sca}$ only, with no scattering matrix at any orientation.

\paragraph{Birefringence comes free with the aligned calculation
(later).} \Ccirc\ is likewise absent from the release, but it does not
need a separate computation. In the forward-scattering amplitude the
difference $S_{\theta\theta} - S_{\phi\phi}$ carries dichroism in its
imaginary part and birefringence in its real part, so a fixed-orientation
amplitude run yields \Cpol\ and \Ccirc\ from one output. Whoever does
the aligned scattering matrix gets \Ccirc\ at no additional cost, which
is why the two belong together whenever they are done. They can be
deferred for now: the 2018 model had no birefringence either, and the
observable being updated is linear polarization.

\paragraph{The averaging conventions do not match, but not yet
critically.} SEDust builds \Cabs\ from the exact orientation average of
its own T-matrix route and \Cpol\ from the MPFA combination of
Eq.~(\ref{eq:cpol}). The two conventions disagree by a median of
$0.022\%$ in the far infrared and $0.080\%$ in the optical, and by a
median of $0.21\%$ with a maximum of $9.53\%$ in the ultraviolet. Across
the optical the $0.080\%$ figure is far below every other uncertainty in
the problem, so this obstructs neither the single-band update nor an
extension to two or three optical bands. It becomes an obstacle as soon
as polarized extinction is quoted in the ultraviolet, where the
polarization curve is measured. Note that
resolving it changes \Cabs, and therefore breaks the byte-identical
regression of the existing SEDs; that is an engineering consequence to
plan around, not a reason to leave two conventions in one model
permanently.

For v1.20 the decision was taken explicitly to defer it. The scope of
the update lies wholly inside the range where the disagreement is
$0.080\%$ or less, and unifying the convention would spend the
regression baseline against which everything else in the release was
validated, in exchange for accuracy nothing in the current plan needs.
The deferral is a statement about ordering, not about correctness: two
orientation averages in one model is still wrong, and the first
ultraviolet polarized-extinction result is the point at which it must be
fixed rather than postponed again.

\subsection{Recommended order}

For the update as scoped, one band first and possibly two or three
optical bands after, the realistic sequence is:

\begin{enumerate}
\item Compute the $V$-band dichroic cross section from physics and
      retire the empirical formula, using the \code{Cpol\_ext} that
      stage~1 already produces.
\item Test the sensitivity to the \falign\ functional form, step versus
      smooth.
\item Bring the random-orientation scattering matrix into the transfer.
      The supply side of this step is complete: \code{run\_scatmat.x}
      already provides $Z(\theta)$ in five optical bands, and widening
      that set costs another run rather than another method. What
      remains is consuming it.
\item Only then the aligned scattering matrix, \Ccirc, the averaging
      convention, and any extension beyond the optical.
\end{enumerate}

The point worth emphasizing is that step~1 on its own is already the
substantive advance over Seon (2018): it replaces a fitted power law
that has no wavelength argument and no alignment physics with a cross
section computed from the grain model, which is exactly the defect
identified in \S\ref{sec:motivation}. Step~1 needs no new T-matrix run,
because \code{Cpol\_ext} exists and has been validated against the HD23
release. Step~3 is placed after step~2 because it is the point at which
the comparison against the release is partly lost, the release
containing no scattering matrix; entering locally computed territory is
better done once the cheaper questions have been answered.

\subsection{What the Hoang group assumes}
\label{sec:hoangvalues}

Since the alignment inputs are assumptions rather than derived
quantities, it is useful to record the values in current use by the
group that produces most of this literature. Table~\ref{tab:hoang}
collects them. Most of these carry over to a single-wavelength
calculation unchanged: \fhighJ, $Q_X$, $Q_J$, the iron-inclusion
parameters and $F_{\rm turb}$ are all wavelength-independent alignment
parameters, so they enter the $V$-band Rayleigh reduction factor exactly
as they enter a submillimeter one. Only the axial ratio and the \falign\
shape have consequences that differ from band to band.

\begin{table}[!ht]
\centering
\caption{Parameter values adopted in the Hoang-group modeling papers.
Where a paper scans a parameter, the scanned range is given.}
\label{tab:hoang}
\begin{tabularx}{\textwidth}{l l X}
\toprule
Quantity & Value & Source and remark \\
\midrule
$Q_X^{\rm high\text{-}J}$, $Q_X^{\rm low\text{-}J}$ &
  $1$, $\simeq 0.3$ &
  \code{2024ApJ...965..183H} Eq.~12; standard case, suprathermal
  rotation taken to give perfect internal alignment \\
same, slow internal relaxation &
  $0.15$, $0.05$ &
  \code{2025A\&A...703A.192T}; wrong internal alignment enters as
  $-0.1$ \\
$Q_J$ (external alignment) &
  $1$ in both states &
  \code{2024ApJ...965..183H}; assumed perfect, so the whole reduction of
  $R$ comes from internal alignment. Stated by the authors as optimistic \\
\fhighJ &
  $0.25$, $0.5$, $1.0$ &
  \code{2023MNRAS.520.3788G} Eq.~31, a step function in
  $\delta_{\rm mag}$; used as model labels for pure RAT, intermediate
  and full MRAT. A footnote there gives $0.25$--$0.7$ for compact
  irregular grains, making $1.0$ an upper scenario \\
$f^{\rm fast}_{\rm high\text{-}J}$, attainable maximum &
  $\sim 0.45$ (SPM), $\sim 0.22$ (PM) &
  \code{2025ApJ...994..115H}; the value $1.0$ used by POLARIS is
  therefore optimistic (\S\ref{sec:mrat}) \\
$N_{\rm cl}$, iron atoms in a cluster &
  $\sim 5$--$10^3$ &
  \code{2025ApJ...980..105G}; reproducing ALMA data needs
  $10^2$--$10^4$ together with $a \geq 10\um$ and inelastic relaxation \\
$\phi_{\rm sp}$, cluster volume filling &
  $0.1$ &
  \code{2025ApJ...980..105G}, with the iron fraction $f_p$ \\
Axial ratio &
  $1.4$ default; $1.6$, $2.0$, $3.0$ scanned &
  \code{2023ApJ...948...55H} for the default; reproducing Pipe-109
  needs $1.4$ at low \Av\ and $2$--$3$ above $\Av = 14$
  (\code{2025A\&A...703A.192T}) \\
$F_{\rm turb}$ &
  $0.98$--$0.91$ &
  $F_{\rm turb} = \frac12[3\langle\cos^2\Delta\theta\rangle - 1]$;
  Pipe-109 values, \code{2025A\&A...703A.192T} \\
$u_{\rm MMP}$ &
  $8.64\times10^{-13}$ erg cm$^{-3}$ &
  normalization of $U = u_{\rm rad}/u_{\rm MMP}$ \\
$T_d$ &
  $16.4\,{\rm K}\,(a/0.1\um)^{-1/15}U^{1/6}$ &
  equilibrium temperature only \\
\falign\ shape &
  step, or
  $f_{\max}[1 - e^{-(0.5a/\aalign)^3}]$ &
  step in POLARIS; the smooth form is \code{2020ApJ...896...44L}, used
  by GRADE-POL on its semi-analytic path \\
\bottomrule
\end{tabularx}
\end{table}

Four features of this table are worth drawing out, because each is a
place where our own model can differ deliberately rather than by
accident.

\paragraph{Negative $Q_X$ is a physical statement, not a bookkeeping
sign.} In the conservative slow-relaxation model the wrong internal
alignment state carries $Q_X = -0.1$, and a negative Rayleigh reduction
factor means the polarization direction flips by $90^\circ$. This is the
origin of the observed pattern in which $P$ is parallel to $\mathbf{B}$
at millimeter wavelengths and perpendicular to it in the submillimeter.
An implementation that clamps $R$ to non-negative values will silently
delete this effect.

\paragraph{External alignment is assumed, not computed.} With $Q_J = 1$
in both the low-$J$ and high-$J$ states, the entire departure of $R$
from unity in these papers comes from internal alignment. The papers say
so; it is an acknowledged optimistic approximation rather than a hidden
one. It also means that quoted values of $R$ are upper bounds.

\paragraph{The radiation field normalization is used loosely.} The group
writes $U = u_{\rm rad}/u_{\rm MMP}$ with
$u_{\rm MMP} = 8.64\times10^{-13}$ erg cm$^{-3}$ throughout, without
distinguishing the $1.0\times10^{-13}$ and $1.65\times10^{-13}$ dilution
factors of the Mathis field and without mentioning the CMB. A survey of
eleven papers from the group found zero that make either distinction.
For SEDust this is a place where the existing radiation-field handling
is already more careful, and the difference should be stated rather than
absorbed.

\paragraph{The temperature treatment is equilibrium throughout.} No
paper in that survey of eleven cites an astrodust heat capacity, and at
least one (\code{2026ApJ..1005..195T}) uses astrodust cross sections
together with a silicate-based temperature scaling. Tram et al.\ (2025)
state explicitly that the temperature distribution was discarded. As
noted in \S\ref{sec:stoch}, SEDust already solves $P(T)$, so the correct
treatment costs one accumulator here and is a genuine improvement over
the reference implementations.

\subsection{A working prescription}

Given those assumptions, the practical question is what to compute and
in what order. The following is written against the current state of our
own code rather than against POLARIS.

\paragraph{1. Start with a fixed \falign.} Adopt the HD23 fit of
Eq.~(\ref{eq:falign}) with $\aalign = 0.0749\um$, $\alpha = 1.80$,
$f_{\max} = 1$, and do not compute $R$ at all. This is already
implemented and verified (\S\ref{sec:codes}). The first optical step is
then not to compute alignment but to swap in the smooth form of
Table~\ref{tab:hoang} and measure how much the polarization curve moves.
If it moves little, the assumed form is not the limiting uncertainty; if
it moves a lot, that settles the priority.

\paragraph{2. Replace the empirical dichroism at $V$.} Use
\code{Cpol\_ext} in Eq.~(\ref{eq:tanh}) in place of
Eq.~(\ref{eq:whittet}), with the alignment weight applied as in
Eq.~(\ref{eq:cpolR}). This is the deliverable: the polarization is then
set by the grain model and the field geometry rather than by a fitted
exponent, and the accumulation becomes multiplicative as
\S\ref{sec:motivation} requires. No new optical data is needed.

\paragraph{3. Only then compute alignment.} Computing $R$ rather than
assuming it needs the analytic chain of \S\ref{sec:alignment},
Eqs.~(\ref{eq:aalign}) and (\ref{eq:avdep}), evaluated on the local
radiation field, which in turn requires the vectorial tally of
\S\ref{sec:vecfield}. \fhighJ\ follows from $\delta_{\rm mag}$, either
through the step function of Giang et al.\ (2023) or the revised
criterion of Eq.~(\ref{eq:deltacri}). This is stage~4 of
Table~\ref{tab:stages} and should not be attempted before steps~1
and~2 have shown what the fixed-\falign\ model cannot do.

\paragraph{4. Climb the validation ladder in order.} Reproduce the
released polarized extinction (done, \S\ref{sec:worklist}); then the
Serkowski $\lambda_{\max}$ and $K$; then $p_V/E(B-V)$ against the
$\geq 13\,\%$/mag requirement of \S\ref{sec:validation}, not the
classical $9\,\%$/mag, which a model that under-polarizes will pass;
then $R_{S/V}$; then the HAWC+ and Planck comparisons. The ordering
matters because each rung uses only quantities the previous rungs have
already validated.

\subsection{A constraint that applies only to a multi-wavelength
extension}

One result deserves to be recorded here even though it is out of scope
for the single-wavelength update, because it will be the first
difficulty encountered as soon as the scope widens. The Pipe-109
analysis found that no single $a_{\max}$ fits all bands:
$0.19$--$0.25\um$ in $R$, $0.25$--$0.35\um$ in $H$, and $1$--$3\um$ at
$850\um$, and that the axial ratio has to grow with \Av, from $1.4$ at
low extinction to $2$--$3$ above $\Av = 14$
(\code{2025A\&A...703A.192T}). Fitting optical, near-infrared and
submillimeter polarization at once therefore strains the size
distribution and the grain shape before it strains the alignment
physics. Updating Seon (2018) at $V$ alone does not run into this: a
single band constrains one combination of size distribution and shape,
and no tension can arise between bands that are not being fitted. The
constraint is listed so that it is not rediscovered as a surprise, and
so that it is not mistaken for a present obstacle.

\paragraph{An honest limit.} Following this prescription does not close
the gap in polarization amplitude. Our intrinsic polarization fraction
tops out at $19.15\%$ at $850\um$ against the Planck maximum of $22\%$
(Table~\ref{tab:validation}). That shortfall is a property of the
$b/a = 1.4$ oblate spheroid and will not be removed by better alignment
physics or a better solver. It is the same pressure that leads the Hoang
group to raise the axial ratio to $2$--$3$ when fitting real sight
lines, and it belongs with the open shape question of
\S\ref{sec:unverified}.

%====================================================================
\section{Validation targets}
\label{sec:validation}

\begin{table}[!ht]
\centering
\caption{Acceptance targets for a polarized dust model.}
\label{tab:validation}
\begin{tabularx}{\textwidth}{l l X}
\toprule
Quantity & Value & Source \\
\midrule
$\lambda_{\max}$ & $0.55$\um & \code{2009ApJ...696....1D} \\
Serkowski $K$ & $-0.01 + 1.66\,\lambda_{\max}$ & \code{1992ApJ...386..562W} \\
$p_{\max}$ at 353 GHz & $22.0^{+3.5}_{-1.4}\ \%$ & \code{2020A\&A...641A..12P} \\
$R_{S/V}$ & $4.2 \pm 0.2$ & \code{2015A\&A...576A.106P} \\
$R_{P/p}$ & $5.4 \pm 0.2$ MJy/sr & \code{2015A\&A...576A.106P} \\
Polarization spectrum, 250\um--3\,mm & flat to $\lesssim 10\%$ & HD23 \\
$p_V / E(B-V)$ & $\geq 13\ \%$/mag & \code{2020A\&A...641A..12P},
  \code{2019A\&A...624L...8P} \\
\bottomrule
\end{tabularx}
\end{table}

\paragraph{Correction to the classical limit.} The last row deserves
emphasis, because using the traditional value will let a bad model pass.
The classical upper bound $p_V/E(B-V) \leq 9\ \%$/mag is obsolete; a
model that under-polarizes will satisfy it comfortably. The post-Planck
requirement is $\geq 13\ \%$/mag, and the HD23 target value is
$0.13$/mag.

\paragraph{Values already computed here.} Three numbers were obtained
during this audit and are available as immediate checks:
\begin{itemize}
\item $Q_{\rm pol}/Q_{\rm abs} = 0.2790$ in the Rayleigh limit for
      $b/a = 1.4$, independent of size and wavelength. This is a pure
      geometry check and should be reproduced exactly.
\item $p_{\rm submm} = 19.5\%$ with \falign\ applied, against the Planck
      value of $22.0\%$.
\item $R_{S/V} = 3.85$ under a single-temperature approximation, against
      $4.2$ observed and $4.31$ as the HD23 target.
\end{itemize}

The third of these is the most informative, and it points to a natural
first task. The $11\%$ shortfall in $R_{S/V}$ is controlled by the
relative emission weighting across grain sizes, and that weighting is
exactly what SEDust's $T(a)$ and $P(T)$ solution provides. Recomputing
$R_{S/V}$ with the full temperature treatment rather than a single $T$ is
therefore both a physics improvement and a sharp test of the new code
path.

%====================================================================
\section{Unverified items}
\label{sec:unverified}

The following are stated as open, not as established. They are listed so
that none of them is silently promoted to fact in later work.

\begin{enumerate}
\item \textbf{$F_{\rm IR}$ coefficient normalization.} Baseline POLARIS
      uses \code{1.40e10} where NewPOLARIS uses \code{6.14e9}. Whether
      this is a difference of unit system or a genuine disagreement was
      not determined. Since $F_{\rm IR}$ enters both
      Eq.~(\ref{eq:aalign}) and Eq.~(\ref{eq:deltacri}), the factor
      matters and must be resolved before either is trusted.

\item \textbf{Origin of the $R_{S/V}$ gap.} Whether $3.85$ versus $4.31$
      is entirely attributable to the single-temperature approximation,
      or partly to a difference in the definition of the ratio, is not
      established.

\item \textbf{Grain shape is an open choice, not a settled one.} Guillet
      et al.\ (2018) favor prolate grains with $a/b = 3$ plus porosity,
      while HD23 use a single-composition oblate spheroid with
      $b/a = 1.4$. These are mutually incompatible solutions to
      overlapping data, and adopting either is a modeling decision that
      should be recorded as such.
\end{enumerate}

%====================================================================
\section{References}

\begin{itemize}\setlength{\itemsep}{2pt}

\item Baes, M., et al.\ 2019, A\&A, 630, A61.
      \code{2019A\&A...630A..61B}

\item Bianchi, S., Ferrara, A., \& Giovanardi, C.\ 1996, ApJ, 465, 127.
      \code{1996ApJ...465..127B}

\item Camps, P., et al.\ 2015, A\&A, 580, A87.
      \code{2015A\&A...580A..87C}

\item Chiar, J.~E., et al.\ 2006, ApJ, 651, 268.
      \code{2006ApJ...651..268C}

\item Draine, B.~T.\ 2003, ApJ, 598, 1017.
      \code{2003ApJ...598.1017D}

\item Draine, B.~T., \& Fraisse, A.~A.\ 2009, ApJ, 696, 1 (DF09).
      \code{2009ApJ...696....1D}

\item Draine, B.~T., \& Hensley, B.~S.\ 2021a, ApJ, 909, 94 (DH21a).
      \code{2021ApJ...909...94D}

\item Draine, B.~T., \& Hensley, B.~S.\ 2021b, ApJ, 919, 65 (DH21b).
      \code{2021ApJ...919...65D}

\item Draine, B.~T., \& Lee, H.~M.\ 1984, ApJ, 285, 89.
      \code{1984ApJ...285...89D}

\item Draine, B.~T., \& Li, A.\ 2007, ApJ, 657, 810 (DL07).
      \code{2007ApJ...657..810D}

\item Giang, N.~C., et al.\ 2025, ApJ, 980, 105.
      \code{2025ApJ...980..105G}

\item Giang, N.~C., Hoang, T., \& Kim, J.-G.\ 2023, MNRAS, 520, 3788
      (\code{2023MNRAS.520.3788G}). NewPOLARIS MRAT and RATD extension,
      Zenodo \code{10.5281/zenodo.14199014}

\item Guillet, V., et al.\ 2018, A\&A, 610, A16.
      \code{2018A\&A...610A..16G}

\item Hensley, B.~S., \& Draine, B.~T.\ 2023, ApJ, 948, 55 (HD23).
      \code{2023ApJ...948...55H}

\item Hoang, T.\ 2025, ApJ, 994, 115 (\code{2025ApJ...994..115H}).

\item Hoang, T., et al.\ 2024, ApJ, 965, 183.
      \code{2024ApJ...965..183H}

\item Hoang, T., Tram, L.~N., \& Lee, H.\ 2021, ApJ, 908, 218.
      \code{2021ApJ...908..218H}

\item Jones, T.~J., et al.\ 2020, AJ, 160, 167.
      \code{2020AJ....160..167J}

\item Kim, J., et al.\ 2023, AJ, 165, 223.
      \code{2023AJ....165..223K}

\item Lazarian, A., \& Hoang, T.\ 2019, ApJ, 883, 122.
      \code{2019ApJ...883..122L}

\item Lee, H., et al.\ 2020, ApJ, 896, 44.
      \code{2020ApJ...896...44L}

\item Peest, C., Camps, P., Stalevski, M., Baes, M., \& Siebenmorgen, R.\
      2017, A\&A, 601, A92. \code{2017A\&A...601A..92P}

\item Peest, C., et al.\ 2023, A\&A, 673, A112.
      \code{2023A\&A...673A.112P}

\item Planck Collaboration Int.\ XXI 2015, A\&A, 576, A106.
      \code{2015A\&A...576A.106P}

\item Planck Collaboration 2019, A\&A, 624, L8.
      \code{2019A\&A...624L...8P}

\item Planck Collaboration 2020, A\&A, 641, A12.
      \code{2020A\&A...641A..12P}

\item Reissl, S., Wolf, S., \& Brauer, R.\ 2016, A\&A, 593, A87
      (POLARIS~I). \code{2016A\&A...593A..87R}

\item Seon, K.-I.\ 2018, ApJ, 862, 87. \code{2018ApJ...862...87S};
      arXiv:1806.06101

\item Serkowski, K., Mathewson, D.~S., \& Ford, V.~L.\ 1975, ApJ, 196,
      261. \code{1975ApJ...196..261S}

\item Tram, L.~N., Hoang, T., \& Lazarian, A.\ 2025, A\&A, 703, A192
      (\code{2025A\&A...703A.192T}). GRADE-POL~I.

\item Truong, B., et al.\ 2025, ApJ, 981, 83.
      \code{2025ApJ...981...83T}

\item Truong, B., et al.\ 2026, ApJ, 1005, 195.
      \code{2026ApJ..1005..195T}

\item Whittet, D.~C.~B., et al.\ 1992, ApJ, 386, 562.
      \code{1992ApJ...386..562W}

\item Whittet, D.~C.~B., et al.\ 2008, ApJ, 674, 304.
      \code{2008ApJ...674..304W}

\item Wood, K.\ 1997, ApJ, 477, L25. \code{1997ApJ...477L..25W}

\item Wood, K., \& Jones, T.~J.\ 1997, AJ, 114, 1405.
      \code{1997AJ....114.1405W}

\item Zubko, V., Dwek, E., \& Arendt, R.~G.\ 2004, ApJS, 152, 211.
      \code{2004ApJS..152..211Z}

\end{itemize}

\end{document}
